\documentclass[11pt,a4paper]{article}

\usepackage[utf8]{inputenc}
\usepackage[T1]{fontenc}
\usepackage[margin=2.5cm]{geometry}
\usepackage{graphicx}
\usepackage{booktabs}
\usepackage{amsmath}
\usepackage{listings}
\usepackage{xcolor}
\usepackage{float}
\usepackage[hidelinks]{hyperref}

\lstset{
  basicstyle=\ttfamily\small,
  breaklines=true,
  frame=single,
  backgroundcolor=\color{gray!8},
  columns=fullflexible
}

\title{Electricity Consumption Forecasting with Recurrent Architectures\\
\large Revised Single-Step Forecasting Experiment}
\author{Erdem YILMAZ \\
Adana Alparslan T\"urke\c{s} Science and Technology University \\
Department of Artificial Intelligence Engineering}
\date{\today}

\begin{document}

\maketitle

\section{GitHub Repository}

\begin{itemize}
  \item \url{https://github.com/eerdemYlmzz/electricity-consumption-forecasting}
\end{itemize}


\section{Revision Compliance}

This report supersedes the previous progress report. The table below maps each item of the
revision request to the action taken and to the file where it can be verified.

\begin{table}[H]
\centering
\small
\begin{tabular}{p{0.30\textwidth}p{0.44\textwidth}p{0.18\textwidth}}
\toprule
\textbf{Revision item} & \textbf{Action} & \textbf{Location} \\
\midrule
1. Redefine the prediction problem &
Task restated as one-hour history $\rightarrow$ next 10-minute step;
\texttt{SEQUENCE\_LENGTH} $=6$, \texttt{FORECAST\_HORIZON} $=1$ &
Section~\ref{sec:problem} \\
\addlinespace
1. Correct the \texttt{\# 60 stride} comment &
Corrected: 60 was the window length, the stride is 1 &
\texttt{lstm-model.ipynb} \\
\addlinespace
1. No shuffling; split before windowing &
Already in place; retained &
Section~\ref{sec:split} \\
\addlinespace
2. Apply normalization &
\texttt{StandardScaler} fitted on the training partition only; metrics computed after
\texttt{inverse\_transform} &
Section~\ref{sec:scaling} \\
\addlinespace
3. Add simple reference models &
Three added: persistence, seasonal naive, moving average &
Section~\ref{sec:references} \\
\addlinespace
4. Simple, explainable architecture &
Main model is hidden $=32$, one layer (4\,513 parameters); a larger model is reported
separately with validation evidence &
Section~\ref{sec:architecture} \\
\addlinespace
4. Remove unused variables &
\texttt{self.tanh}, \texttt{self.relu} and \texttt{self.input\_size} removed &
\texttt{lstm-model.ipynb} \\
\addlinespace
5. Describe the prediction mode correctly &
Stated as sliding-window one-step-ahead on observed values; explicitly not a long-horizon
forecast &
Section~\ref{sec:problem} \\
\addlinespace
6. Improve the figures &
Five figures produced, all with real date-time axes &
Section~\ref{sec:figures} \\
\addlinespace
6. Do not label the axis kWh &
The unit is not documented in the source; all axes and tables now read ``Consumption'' &
Section~\ref{sec:unit} \\
\bottomrule
\end{tabular}
\caption{Revision items and the corresponding changes.}
\end{table}

\noindent
\textbf{Note on item 3.} In the revision document the sentence ``the results should be compared
with at least two simple methods'' is followed by an empty bullet list, so the intended methods
were not specified. Three standard references were selected and are defined explicitly in
Section~\ref{sec:references}.


\section{Problem Definition}
\label{sec:problem}

The series is sampled every 10 minutes, so one hour of history corresponds to six samples. The
task is:

\begin{quote}
Using the consumption values of the past hour, predict the consumption value of the next
10-minute step.
\end{quote}

\begin{table}[H]
\centering
\begin{tabular}{lcl}
\toprule
\textbf{Parameter} & \textbf{Value} & \textbf{Meaning} \\
\midrule
\texttt{SEQUENCE\_LENGTH}  & 6 & window length: $6 \times 10$ minutes $=$ 1 hour of history \\
\texttt{FORECAST\_HORIZON} & 1 & one step ahead, i.e.\ 10 minutes \\
Stride                     & 1 & the window advances by one sample \\
\bottomrule
\end{tabular}
\caption{Task parameters. The dataset has no one-minute observations, so a one-minute-ahead
problem cannot be posed on this data; the smallest available horizon is 10 minutes.}
\end{table}

\subsection{How the predictions are produced}

This is \textbf{sliding-window one-step-ahead prediction}. Every prediction is computed from six
\emph{observed} measurements. The model is never fed its own output as input, and the test period
is not generated in one pass.

\begin{quote}
The model predicts the consumption 10 minutes ahead \emph{given that the actual measurements of
the previous hour are known}.
\end{quote}

\noindent
The results must therefore not be read as a long-horizon forecast. Free-running (autoregressive)
multi-step forecasting is a different experiment and has not been performed.

\subsection{Correction to the previous configuration}

The earlier code contained:

\begin{lstlisting}[language=Python]
SEQUENCE_LENGTH = 60  # 60 stride
\end{lstlisting}

\noindent
The comment was wrong: 60 was the \emph{window length} in samples, not the stride. The stride in
that implementation was 1, as it is here. The comment has been corrected in the notebook, and the
revised experiment uses a window of 6.


\section{Dataset}

\subsection{Source}

The dataset is the \emph{Power Consumption of Tetouan City} dataset.

\begin{itemize}
  \item Kaggle: \url{https://www.kaggle.com/datasets/fedesoriano/electric-power-consumption}
  \item UCI Machine Learning Repository:
        \url{https://archive.ics.uci.edu/dataset/849/power+consumption+of+tetouan+city}
\end{itemize}

\subsection{Related Publication}

\begin{quote}
A. Salam and A. E. Hibaoui, ``Comparison of Machine Learning Algorithms for the Power
Consumption Prediction: --- Case Study of Tetouan city ---,'' \emph{2018 6th International
Renewable and Sustainable Energy Conference (IRSEC)}, Rabat, Morocco, 2018, pp. 1--5,
doi: 10.1109/IRSEC.2018.8703007.
\end{quote}

\subsection{Original Features}

Temperature, Humidity, Wind Speed, General Diffuse Flows, Diffuse Flows, and Power Consumption for
Zones 1, 2 and 3. Only \textbf{Zone 1} consumption is used in this study.

\subsection{Unit of the Target Variable}
\label{sec:unit}

The source dataset does not document the physical unit of the consumption columns. The values were
divided by 1000 during preprocessing (the raw Kaggle value $34055.6962$ corresponds to $34.055696$
in the processed file), which changes the magnitude but does not establish the unit. Accordingly,
\textbf{no physical unit is claimed anywhere in this report}: all axis labels and table headings
read ``Consumption''. Confirming the unit against the dataset documentation remains an open item.

\subsection{Timestamp Correction}

The \texttt{datetime} column of \texttt{data/data\_processed.csv} does not match the source data.
Preprocessing kept only the hour and minute,

\begin{lstlisting}[language=Python]
df["time"] = df["time"].apply(lambda x: pd.to_datetime(x).strftime("%H:%M"))
\end{lstlisting}

\noindent
and the date component was subsequently regenerated, producing a synthetic range of
2025-07-01 to 2026-06-29. The true range of the source dataset is 2017-01-01 to 2017-12-30.

Because the revision requires real dates and times on the plot axes, plotting synthetic dates
would have been worse than plotting sample indices. The true index is therefore reconstructed in
\texttt{lstm\_forecast.py}. This is exact rather than approximate: the file contains 52\,416 rows
at a 10-minute stride, which is exactly 364 days, so

\begin{lstlisting}[language=Python]
pd.date_range("2017-01-01 00:00", periods=52416, freq="10min")
\end{lstlisting}

\noindent
reproduces 2017-01-01 00:00 to 2017-12-30 23:50. The ordering and spacing of the samples were never
affected by the preprocessing error; only the labels were. Regenerating the \texttt{datetime}
column correctly in the preprocessing stage is recommended.

\subsection{Data Used}

\begin{table}[H]
\centering
\begin{tabular}{ll}
\toprule
\textbf{Property} & \textbf{Value} \\
\midrule
Shape & $52416 \times 3$ \\
Sampling interval & 10 minutes \\
True date range & 2017-01-01 00:00 to 2017-12-30 23:50 (364 days) \\
\texttt{consumption} & stored as \texttt{float64}, e.g.\ 34.055696 (unit undocumented) \\
Feature count & 1 (Zone 1 consumption only) \\
Scaling & \texttt{StandardScaler}, fitted on the training partition only \\
\bottomrule
\end{tabular}
\caption{Structure of the data used for training.}
\end{table}


\section{Preprocessing}

\subsection{Chronological Split}
\label{sec:split}

The raw series is split chronologically into 70\,/\,10\,/\,20 partitions \textbf{before} windowing,
and the sliding windows are generated independently inside each partition. Nothing is shuffled.

If the windows were generated first and split afterwards, a window near a boundary would take its
inputs from one partition and its target from the next, leaking training timesteps into validation
and test.

\begin{table}[H]
\centering
\small
\begin{tabular}{lccccc}
\toprule
\textbf{Partition} & \textbf{Raw index} & \textbf{Length} & \textbf{Sequences} & \textbf{First timestamp} & \textbf{Last timestamp} \\
\midrule
Train      & $[0,\ 36691)$     & 36691 & 36685 & 2017-01-01 00:00 & 2017-09-12 19:00 \\
Validation & $[36691,\ 41932)$ & 5241  & 5235  & 2017-09-12 19:10 & 2017-10-19 04:30 \\
Test       & $[41932,\ 52416)$ & 10484 & 10478 & 2017-10-19 04:40 & 2017-12-30 23:50 \\
\bottomrule
\end{tabular}
\caption{Leakage-free chronological split. Each partition loses
\texttt{SEQUENCE\_LENGTH} $=6$ samples to windowing.}
\end{table}

\subsection{Normalization}
\label{sec:scaling}

A \texttt{StandardScaler} is fitted on the training partition only and then applied to all three
partitions:

\begin{lstlisting}[language=Python]
from sklearn.preprocessing import StandardScaler

scaler = StandardScaler()
scaler.fit(train_values)

train_scaled = scaler.transform(train_values)
val_scaled   = scaler.transform(val_values)
test_scaled  = scaler.transform(test_values)
\end{lstlisting}

\noindent
Fitting the scaler on the full series would carry the mean and variance of the validation and test
periods into training and constitute leakage. The fitted statistics are
$\mu = 32.9569$ and $\sigma = 7.2555$, computed on the training partition alone.

Before any metric is computed, both predictions and targets are transformed back:

\begin{lstlisting}[language=Python]
y_pred_original = scaler.inverse_transform(y_pred_scaled)
y_true_original = scaler.inverse_transform(y_true_scaled)
\end{lstlisting}

\noindent
Every MAE, RMSE and MAPE value in this report is therefore expressed in the original consumption
scale, and is directly comparable to the reference models.


\section{Reference Models}
\label{sec:references}

Three training-free references are evaluated on exactly the same targets as the network. All three
use only past values, so none of them leak.

\begin{table}[H]
\centering
\begin{tabular}{lll}
\toprule
\textbf{Reference} & \textbf{Prediction} & \textbf{Interpretation} \\
\midrule
Persistence     & $\hat{y}_t = y_{t-1}$                        & repeat the last observed value \\
Seasonal naive  & $\hat{y}_t = y_{t-144}$                      & same time of day, previous day \\
Moving average  & $\hat{y}_t = \frac{1}{6}\sum_{k=1}^{6} y_{t-k}$ & mean of the input window \\
\bottomrule
\end{tabular}
\caption{Reference models. $144 = 24\,\mathrm{h} / 10\,\mathrm{min}$ is one full day of samples.}
\end{table}

\begin{table}[H]
\centering
\begin{tabular}{lccccc}
\toprule
\textbf{Reference} & \textbf{MSE} & \textbf{RMSE} & \textbf{MAE} & \textbf{MAPE} & \textbf{R2} \\
\midrule
Persistence    & 0.3025 & 0.5500 & 0.3696 & 1.27\% & 0.9920 \\
Seasonal naive & 2.2303 & 1.4934 & 1.0058 & 3.47\% & 0.9414 \\
Moving average & 2.3903 & 1.5461 & 1.1002 & 3.74\% & 0.9372 \\
\bottomrule
\end{tabular}
\caption{Reference models on the test partition, original consumption scale.}
\end{table}

\noindent
Persistence is by a wide margin the strongest of the three at this horizon: over 10 minutes the
last observed value carries far more information than either the previous day's value at the same
time or the mean of the past hour. The moving average is the weakest because averaging six samples
lags the signal during the steep morning and evening transitions. \textbf{Persistence is therefore
the threshold the trained model has to beat}; comparing against the two weaker references alone
would overstate the model's value.


\section{Model Architecture}
\label{sec:architecture}

The input is a single feature over six timesteps, so the revision specifies a small starting
model. That configuration is the main model of this report.

\begin{table}[H]
\centering
\begin{tabular}{lcc}
\toprule
\textbf{Parameter} & \textbf{Main model} & \textbf{Capacity check} \\
\midrule
\texttt{input\_size}   & 1 & 1 \\
\texttt{hidden\_size}  & 32 & 64 \\
\texttt{num\_layers}   & 1 & 2 \\
\texttt{output\_size}  & 1 & 1 \\
Trainable parameters   & 4\,513 & 50\,497 \\
\bottomrule
\end{tabular}
\caption{Architectures. Both are trained under an identical protocol: same seed, split, scaler,
optimizer, epoch budget and stopping rule.}
\end{table}

\noindent
The unused attributes of the earlier implementation --- \texttt{self.tanh}, \texttt{self.relu} and
\texttt{self.input\_size} --- were defined in the constructor but never referenced in
\texttt{forward()}. They have been removed; the class now defines only the LSTM and the output
layer.


\section{Training Setup}

\begin{table}[H]
\centering
\begin{tabular}{ll}
\toprule
\textbf{Parameter} & \textbf{Value} \\
\midrule
Loss function & MSELoss \\
Optimizer & Adam \\
Learning rate & $1 \times 10^{-3}$ \\
Batch size & 64 \\
Shuffle & False \\
Maximum epochs & 50 \\
Early stopping patience & 8 epochs on validation loss \\
Seed & 42 (\texttt{random}, \texttt{numpy}, \texttt{torch}) \\
\bottomrule
\end{tabular}
\caption{Training hyperparameters.}
\end{table}

\noindent
The weights of the best validation epoch are restored before the test partition is evaluated. The
seed is fixed so that the two architectures differ only in capacity and not in initialisation.


\section{Results}

\begin{table}[H]
\centering
\begin{tabular}{lccccc}
\toprule
\textbf{Model} & \textbf{MSE} & \textbf{RMSE} & \textbf{MAE} & \textbf{MAPE} & \textbf{R2} \\
\midrule
Persistence                   & 0.3025 & 0.5500 & 0.3696 & 1.27\% & 0.9920 \\
Seasonal naive                & 2.2303 & 1.4934 & 1.0058 & 3.47\% & 0.9414 \\
Moving average                & 2.3903 & 1.5461 & 1.1002 & 3.74\% & 0.9372 \\
\midrule
LSTM, hidden 32, 1 layer      & 0.1663 & 0.4078 & 0.2830 & 0.99\% & 0.9956 \\
LSTM, hidden 64, 2 layers     & 0.1307 & 0.3616 & 0.2556 & 0.90\% & 0.9966 \\
\bottomrule
\end{tabular}
\caption{Test partition, original consumption scale. The test partition was evaluated once, after
training and model selection were complete.}
\end{table}

\begin{table}[H]
\centering
\begin{tabular}{lccccc}
\toprule
\textbf{Model} & \textbf{MSE} & \textbf{RMSE} & \textbf{MAE} & \textbf{MAPE} & \textbf{R2} \\
\midrule
Persistence                   & 0.3814 & 0.6176 & 0.4162 & 1.23\% & 0.9910 \\
Seasonal naive                & 2.8999 & 1.7029 & 1.2041 & 3.67\% & 0.9315 \\
Moving average                & 3.2015 & 1.7893 & 1.2620 & 3.67\% & 0.9244 \\
\midrule
LSTM, hidden 32, 1 layer      & 0.1517 & 0.3895 & 0.2610 & 0.79\% & 0.9964 \\
LSTM, hidden 64, 2 layers     & 0.1192 & 0.3452 & 0.2375 & 0.73\% & 0.9972 \\
\bottomrule
\end{tabular}
\caption{Validation partition at the selected epoch. These values are optimistically biased,
because the epoch was chosen by minimising validation loss.}
\end{table}

\begin{table}[H]
\centering
\begin{tabular}{lcccc}
\toprule
\textbf{Model} & \textbf{Parameters} & \textbf{Epochs run} & \textbf{Best epoch} & \textbf{Stopped early} \\
\midrule
LSTM, hidden 32, 1 layer  & 4\,513  & 20 & 12 & yes \\
LSTM, hidden 64, 2 layers & 50\,497 & 50 & 47 & no \\
\bottomrule
\end{tabular}
\caption{Training behaviour (maximum 50 epochs, patience 8).}
\end{table}

\subsection{Justification for the Larger Model}

The revision requires that any model larger than the mandated baseline be justified with
validation results. The comparison is as follows.

\begin{table}[H]
\centering
\begin{tabular}{lccc}
\toprule
\textbf{Quantity} & \textbf{hidden 32, 1 layer} & \textbf{hidden 64, 2 layers} & \textbf{Change} \\
\midrule
Validation RMSE & 0.3895 & 0.3452 & $-11.4\%$ \\
Test RMSE       & 0.4078 & 0.3616 & $-11.3\%$ \\
Parameters      & 4\,513 & 50\,497 & $\times 11.2$ \\
\bottomrule
\end{tabular}
\caption{Capacity comparison. The ordering established on validation is confirmed on test.}
\end{table}

\noindent
The larger model is better on validation by 11.4\%, and the ordering chosen on validation is
confirmed on the test partition, which is the meaningful check. Two qualifications apply:

\begin{itemize}
  \item The larger model did \textbf{not} trigger early stopping: it consumed the full 50-epoch
        budget and its best epoch was 47, so it was still improving when training ended. The
        reported gap is therefore a lower bound on its advantage, not a converged measurement.
  \item The improvement costs 11.2 times more parameters for an 11.3\% reduction in test RMSE.
        Whether that trade is worthwhile depends on the purpose. For the architecture comparison
        planned next (GRU, BiGRU, attention, TCN), the smaller model is the better control because
        it converges in 20 epochs and its results are reproducible within a short budget.
\end{itemize}

\noindent
The main model of this report is therefore the mandated hidden $=32$, one-layer configuration, and
the larger configuration is reported as evidence rather than adopted.


\section{Figures}
\label{sec:figures}

All figures below use the main model (hidden 32, one layer). The horizontal axis shows real dates
and times, and the vertical axis is labelled ``Consumption'' because the unit is undocumented.

\begin{figure}[H]
\centering
\includegraphics[width=\textwidth]{revised_loss_curves.png}
\caption{Training and validation loss per epoch for both architectures. The dotted line marks the
epoch whose weights were restored. Losses are in the standardized scale in which the network is
trained. In the left panel the validation loss turns upward after epoch 12 while the training loss
keeps falling, which is what triggers early stopping; the rise is small (0.002882 to 0.002899).
In the right panel the larger model is still descending at epoch 50.}
\label{fig:loss}
\end{figure}

\begin{figure}[H]
\centering
\includegraphics[width=\textwidth]{revised_test_actual_vs_predicted.png}
\caption{Actual and predicted consumption over the whole test period, 2017-10-19 to 2017-12-30.}
\label{fig:testfull}
\end{figure}

\begin{figure}[H]
\centering
\includegraphics[width=\textwidth]{revised_test_week_zoom.png}
\caption{First week of the test period, 2017-10-19 to 2017-10-26. The daily cycle and the sharp
evening peak are both reproduced; the visible discrepancies are concentrated at the turning
points.}
\label{fig:week}
\end{figure}

\begin{figure}[H]
\centering
\includegraphics[width=\textwidth]{revised_test_errors.png}
\caption{Prediction error over time (top) and the distribution of the errors (bottom). The error
series is centred on zero with no visible trend across the test period; the largest excursions
coincide with rapid transitions in the daily cycle.}
\label{fig:errors}
\end{figure}

\begin{figure}[H]
\centering
\includegraphics[width=0.72\textwidth]{revised_test_scatter.png}
\caption{Predicted against actual consumption on the test partition, with the $y=x$ reference
line. Spread around the line is roughly uniform across the range, indicating no systematic bias at
high or low consumption levels.}
\label{fig:scatter}
\end{figure}


\section{Findings}

\begin{enumerate}
  \item \textbf{Both trained models beat every reference.} Against persistence --- the strongest
        reference --- the mandated model reduces test RMSE by 25.9\% (0.5500 to 0.4078) and the
        larger model by 34.3\% (0.5500 to 0.3616). Against seasonal naive and moving average the
        margin exceeds 70\%.

  \item \textbf{A high $R^2$ is not by itself evidence of a good model on this dataset.}
        Persistence reaches $R^2 = 0.9920$ and $\mathrm{MAPE} = 1.27\%$ without any training,
        because consecutive 10-minute values differ very little relative to their magnitude. Only
        the margin over persistence is informative, which is why the reference models are reported
        alongside every result.

  \item \textbf{Shortening the window from 60 to 6 improved the result.} Under the previous
        configuration the same architecture (hidden 64, 2 layers, standardized) reached a test RMSE
        of 0.4047 with a 60-step window; with a 6-step window it reaches 0.3616, an improvement of
        10.7\%. Persistence scores 0.5507 and 0.5500 on the two target sets respectively, so the
        two experiments are directly comparable. For a 10-minute-ahead prediction the additional 54
        steps of history carried no usable information and made optimisation harder.

  \item \textbf{The seasonal reference is weak at this horizon, and this bounds what the
        experiment can show.} Predicting from the same time on the previous day gives RMSE 1.4934
        against 0.5500 for persistence. At a 10-minute horizon the recent value dominates the daily
        pattern, so the experiment cannot demonstrate that the model has learned daily or weekly
        structure. Demonstrating that requires a longer horizon, where persistence degrades and the
        seasonal reference becomes competitive.

  \item \textbf{Early stopping behaved differently for the two capacities.} The small model reached
        its best validation loss at epoch 12 and stopped at epoch 20; the large model was still
        improving at epoch 47 and exhausted the budget. The overfitting signature is therefore
        visible in the small model and absent in the large one within this budget, which is the
        opposite of the naive expectation and follows from the large model's slower convergence.
        The signature is real but slight: the small model's validation loss rises from 0.002882 at
        epoch 12 to 0.002899 at epoch 19, a change of 0.6\%, while its training loss falls from
        0.003867 to 0.003710 over the same span. Overfitting is not what limits accuracy here.

  \item \textbf{The validation--test gap persists.} The main model reports RMSE 0.3895 on
        validation and 0.4078 on test. Two effects contribute: the epoch was selected by minimising
        validation loss, which biases the validation figure optimistically, and the test partition
        is chronologically later (mid-October to end of December against mid-September to
        mid-October), so seasonal drift is present. Notably the reference models show the opposite
        movement --- persistence improves from 0.6176 on validation to 0.5500 on test --- so the
        test period is intrinsically easier for a trivial predictor while being harder for the
        model. This is worth isolating.
\end{enumerate}


\section{Prior Experiments (Superseded)}
\label{sec:ablation}

The experiments in this section were run under the previous problem definition (60-step window,
no forecast-horizon parameter). They are retained because they are the evidence base for two
decisions carried into the revised setup: applying normalization, and using early stopping instead
of a fixed epoch count. Their numbers are not comparable to Section~8 except where explicitly
stated.

\subsection{Scaling Ablation}

Two runs shared seed, architecture, optimizer, epoch budget and stopping policy, and differed only
in whether standardization was applied.

\begin{table}[H]
\centering
\begin{tabular}{lccccc}
\toprule
\textbf{Model} & \textbf{MSE} & \textbf{RMSE} & \textbf{MAE} & \textbf{MAPE} & \textbf{R2} \\
\midrule
Persistence baseline   & 0.3032 & 0.5507 & 0.3697 & 1.27\% & 0.9920 \\
Run A --- no scaling   & 0.1820 & 0.4266 & 0.3278 & 1.18\% & 0.9952 \\
Run B --- standardized & 0.1638 & 0.4047 & 0.2942 & 1.04\% & 0.9957 \\
\bottomrule
\end{tabular}
\caption{Test metrics under the previous 60-step configuration.}
\end{table}

\begin{figure}[H]
\centering
\includegraphics[width=\textwidth]{ablation_val_rmse.png}
\caption{Validation RMSE per epoch for both ablation runs against the persistence baseline,
previous configuration.}
\label{fig:ablation}
\end{figure}

\subsection{Conclusions Carried Forward}

\begin{itemize}
  \item \textbf{The main effect of scaling is convergence speed, not final quality.} Run B reached
        a validation RMSE of 0.4139 at epoch 2, a level Run A only reached at epoch 12. After both
        were trained to convergence the remaining gap in test RMSE was 5.1\%. An earlier version of
        that experiment used a fixed 10-epoch budget and reported a much larger gap (0.5318 against
        0.4318); that difference was an artefact of stopping before the unscaled run had converged.
        Normalization is applied unconditionally in the revised setup, so this is no longer an
        experimental variable.

  \item \textbf{Why the unscaled run converged slowly.} The network is initialised near zero while
        the targets are centred near 33, and Adam bounds each parameter step at approximately the
        learning rate, so the first epochs are spent translating the output to the correct mean
        rather than learning temporal structure. Standardization removes this offset.

  \item \textbf{A fixed epoch count is not a safe training policy here.} With a 10-epoch budget the
        validation loss never diverged from the training loss and no overfitting was observable.
        With a 50-epoch budget the validation loss reached its minimum at epoch 38 in Run A and
        epoch 24 in Run B and then rose while the training loss continued to fall. The effect was
        mild --- Run A validation loss rose from 0.1123 to 0.1157 between epochs 38 and 46 --- but
        it is the reason early stopping is used throughout.
\end{itemize}


\section{Outputs}

\begin{itemize}
  \item Revised experiment: \texttt{lstm\_forecast.py}. Reference models and both architectures are
        run from this single script.
  \item Figures saved at \texttt{dpi = 300} under \texttt{plots/}:
        \texttt{revised\_loss\_curves.png},
        \texttt{revised\_test\_actual\_vs\_predicted.png},
        \texttt{revised\_test\_week\_zoom.png},
        \texttt{revised\_test\_errors.png},
        \texttt{revised\_test\_scatter.png}.
  \item Metrics, hyperparameters and per-epoch history written to
        \texttt{hyperparameter\_logs.json} under the experiment names
        \texttt{reference\_persistence}, \texttt{reference\_seasonal\_naive},
        \texttt{reference\_moving\_average}, \texttt{lstm\_h32\_l1} and \texttt{lstm\_h64\_l2}.
  \item Models saved to \texttt{model/lstm\_h32\_l1.pth} and \texttt{model/lstm\_h64\_l2.pth}.
  \item Superseded ablation: \texttt{lstm\_ablation.py}, logged under
        \texttt{persistence\_baseline}, \texttt{lstm\_v4\_noscale\_earlystop} and
        \texttt{lstm\_v5\_scaled\_earlystop}.
\end{itemize}


\section{Planned Work}

\begin{enumerate}
  \item \textbf{Architecture comparison} --- GRU, BiGRU, attention, TCN --- under the protocol
        fixed in this report: same split, same scaler, same seed, same stopping rule, and the
        reference models reported alongside every result.
  \item \textbf{Longer horizons.} At \texttt{FORECAST\_HORIZON} $=1$ the persistence reference is
        already at $R^2 = 0.9920$, so all architectures will be compressed into a narrow band and
        the differences between them may not be measurable. Evaluating at 1 hour, 6 hours and 24
        hours ahead would open the margin in which architectural differences can be observed, and
        would make the seasonal reference meaningful.
  \item \textbf{Multivariate input} --- temperature, humidity, wind speed, general diffuse flows,
        diffuse flows, and Zones 2--3.
  \item \textbf{XAI-based feature selection} (SHAP, Integrated Gradients, permutation importance),
        which requires the multivariate input above and is therefore a subsequent step rather than
        a parallel one.
  \item \textbf{Autoregressive multi-step forecasting}, in which the model consumes its own
        predictions. This is a genuinely different experiment from the present one and has not been
        attempted.
\end{enumerate}


\section{Open Questions}

\subsection{Resolved by This Revision}

\begin{itemize}
  \item \textbf{Problem definition.} Fixed as one-hour history to one 10-minute step ahead, with
        the window length and horizon as explicit named parameters.
  \item \textbf{The \texttt{\# 60 stride} comment.} Corrected; the stride was always 1.
  \item \textbf{Normalization.} \texttt{StandardScaler} fitted on the training partition only,
        with metrics computed after inverse transformation.
  \item \textbf{Reference models.} Three added, and reported alongside every result.
  \item \textbf{Model size.} Reduced to hidden 32, one layer; the larger configuration is reported
        with validation evidence rather than assumed.
  \item \textbf{Dead code.} \texttt{self.tanh}, \texttt{self.relu} and \texttt{self.input\_size}
        removed.
  \item \textbf{Figures.} Five figures with real date-time axes.
  \item \textbf{Unit labelling.} No physical unit is claimed; axes read ``Consumption''.
  \item \textbf{Validation split, early stopping, per-epoch validation loss.} In place since the
        previous report and retained.
\end{itemize}

\subsection{Still Open}

\begin{itemize}
  \item \textbf{The unit of the consumption column} is still undocumented and should be confirmed
        against the dataset description.
  \item \textbf{The \texttt{datetime} column in \texttt{data\_processed.csv} is synthetic} and
        should be regenerated correctly in the preprocessing stage; the reconstruction currently
        lives in \texttt{lstm\_forecast.py}.
  \item \textbf{The larger model did not converge} within the 50-epoch budget, so the capacity
        comparison is a lower bound rather than a settled result.
  \item \textbf{Single feature.} Only Zone 1 consumption is used.
  \item \textbf{Single seed.} All runs use seed 42. Repeating each configuration over several seeds
        would separate real differences from initialisation noise; this matters most once
        architectures with small expected differences are compared.
  \item \textbf{Direction of the validation--test gap.} The model degrades from validation to test
        while persistence improves, which has not yet been explained.
  \item \textbf{Attention variants} to be used in the comparison have not been chosen.
\end{itemize}


\section{Discussion Notes}

This section records the questions raised during the analysis and the conclusions reached.

\subsection*{Q1. Are all consumption values stored in 4 bytes?}

No. Both \texttt{data/processed\_powerconsumption.parquet} and \texttt{data/data\_processed.csv}
store \texttt{consumption} as \texttt{float64}, that is 8 bytes per value. The conversion to
\texttt{float32} happens only when the tensor is created for training. The \texttt{datetime}
column is stored as a string rather than a datetime type, so it is re-parsed on every read.

\subsection*{Q2. Is it reasonable to lower the learning rate because the unscaled values are large?}

No, and the reasoning does not apply to this optimizer. Adam's update is
$\Delta\theta = \eta \cdot \hat{m} / (\sqrt{\hat{v}} + \epsilon)$. Multiplying all gradients by a
constant scales $\hat{m}$ and $\sqrt{\hat{v}}$ equally, so the step size is unchanged; Adam is
invariant to the global gradient scale. The argument would be correct for plain SGD, where
$\Delta\theta = \eta \cdot g$ scales directly with the gradient magnitude.

The actual problem is the distance the output must travel: the network output starts near zero and
the targets are centred near 33, while Adam bounds each parameter step at approximately
$\eta = 10^{-3}$. Lowering the learning rate shortens each step and therefore makes the problem
worse. Standardization is the correct remedy, and it is now applied unconditionally.

\subsection*{Q3. Does the model appear to have memorised the data?}

Not in the classical sense. The visually near-perfect overlap in Figure~\ref{fig:week} is largely a
property of the task: a persistence predictor produces a similar-looking curve. A mild overfitting
onset does exist and is visible in Figure~\ref{fig:loss} for the small model after epoch 12, which
is what early stopping handles.

\subsection*{Q4. How can the effect of each individual change be isolated?}

By single-variable ablation with a fixed seed: only one independent variable changes between two
runs while architecture, seed, optimizer, epoch budget and stopping policy remain identical. This
is how the scaling question was settled, and how the capacity comparison in
Section~\ref{sec:architecture} is constructed.

\subsection*{Q5. Can MAPE exceed 100\%?}

Yes. MAPE is $\frac{100}{n}\sum |(y_i - \hat{y}_i)/y_i|$, and each term is the error relative to a
reference value, not a part-of-a-whole ratio. If the true value is 30 and the prediction is 96, the
term is $66/30 = 2.2$, i.e.\ 220\%. MAPE is bounded below by 0 and unbounded above. The values
around 1\% observed here are small because the mean consumption is roughly 33 while the mean
absolute error is roughly 0.28.

\subsection*{Q6. Why is MAPE expressed as a percentage rather than a raw ratio?}

The factor of 100 is a presentation convention and carries no information; 0.0099 and 0.99\% are
the same quantity. The convention matters in practice because libraries disagree:
\texttt{sklearn.metrics.mean\_absolute\_percentage\_error} returns the fraction, not the
percentage. The implementation here multiplies by 100 and therefore reports percentages.

\subsection*{Q7. Does the validation set consist of training data?}

No; the partitions are disjoint. The validation partition was carved out of the block previously
assigned to training, but once separated it is excluded from all gradient updates. Its exact index
range and timestamps are given in Section~\ref{sec:split}.

\subsection*{Q8. What exactly distinguishes the validation set from the test set?}

Both are held out from gradient updates, so that is not the distinction. The distinction is whether
a decision was made using the set. A held-out set gives an unbiased estimate of generalisation
error only if nothing was selected using it. Weights are chosen by minimising the training loss, so
the training loss is biased with respect to the weights. The stopping epoch is chosen by minimising
the validation loss, so the validation loss is biased with respect to that choice. The test set is
used for no selection and therefore remains unbiased. Operationally: the validation set may be
inspected and acted upon freely, the test set is inspected once and never acted upon.

\subsection*{Q9. Where does XAI fit into feature selection?}

XAI methods such as SHAP, Integrated Gradients or permutation importance attribute a prediction to
individual input features. The current model has a single input feature, so there is no feature set
to attribute over and no selection decision to support. XAI becomes applicable once the input is
extended to the remaining variables, which is therefore a prerequisite step rather than a parallel
one.

\end{document}
